% !PRINTED 643
\weberpage{643}

%<<CONT
\gesperrt{eine rationale Funktion $n$ten Grades von $t$ ist, deren Koeffizienten $\overline{A}_1,\allowbreak \overline{A}_2,\allowbreak \ldots,\allowbreak \overline{A}_n$ Funktionale des Körpers $\overline{Z}$ sind.}

Wir definieren die Norm und die Spur des Funktionals wie in § 172:
\weq{7}{N(\overline{\omega}) = \left|\begin{array}{llll}
\overline{x}_{1,1}, & \overline{x}_{1,2}, & \ldots, & \overline{x}_{1,n} \\
\overline{x}_{2,1}, & x_{2,2}, & \ldots, & \overline{x}_{2,n} \\
\cdot & \cdot & \cdots & \cdot \\
\overline{x}_{n,1}, & \overline{x}_{n,2}, & \ldots, & \overline{x}_{n,n}
\end{array}\right|} % [sic] second entry of second row printed without overline, reproduced as scanned
\weq{8}{S(\overline{\omega}) = \overline{x}_{1,1} + \overline{x}_{2,2} + \cdots + \overline{x}_{n,n},}
und die Sätze § 172, (6) und (21) gelten unverändert auch von den Normen und Spuren der Funktionale.

Es ist dann, wie in § 172, (8),
\weq{9}{\varphi(t) = N(t - \overline{\omega}),}
und die \gesperrt{Diskriminante eines Funktionalsystems} ist wie in § 173, (1) erklärt, und wenn das Funktional $\overline{\omega}$ in einem in $\overline{\Omega}$ enthaltenen Körper $e$ten Grades $\overline{\Omega}_1$ enthalten ist, so genügt $\overline{\omega}$ einer Gleichung $e$ten Grades $\varphi_1(t) = 0$. Wie in § 172, (19) wird bewiesen, daß
\[
\varphi(t) = [\varphi_1(t)]^f \quad (n = ef)
\]
eine Potenz von $\varphi_1(t)$ ist.

\begin{center}
§ 179.\\
\gesperrt{Ganze Funktionale des Körpers $\overline{\Omega}$.}
\end{center}

6. \gesperrt{Definition. Ein Funktional $\overline{\omega}$ in $\overline{\Omega}$ heißt ganz, wenn es einer Gleichung:}
\weq{1}{F(\overline{\omega}) = \overline{\omega}^m + \overline{A}_1\overline{\omega}^{m-1} + \cdots + \overline{A}_{m-1}\overline{\omega} + \overline{A}_m = 0}
\gesperrt{genügt, deren Koeffizienten $\overline{A}_1, \ldots, \overline{A}_m$ ganze rationale Funktionale sind.}

Aus dem \gesperrt{Gaussschen} Satz (Bd. I, § 2) folgt, da man nach § 177 die Primfaktoren in $\overline{Z}$ kennt, daß, wenn $F(t)$ in $\overline{Z}$ reduzibel ist, auch jede in $F(t)$ aufgehende Funktion ganze Koeffizienten hat, insbesondere also auch die Gleichung niedrigsten Grades $\varphi(\overline{\omega}) = 0$, der $\overline{\omega}$ genügt.

Die notwendige und hinreichende Bedingung für ein ganzes Funktional $\overline{\omega}$ ist also die, daß
\[
N(t - \overline{\omega})
\]
bei Ordnung nach Potenzen von $t$ ganze rationale Funktionale zu Koeffizienten hat.
